\documentclass[11pt]{article}
\usepackage{../../estilo/vanguarda44}

\renewcommand{\contentsname}{Sumário}

\title{\bfseries\color{vinho}\Huge VANGUARDA 44: ARDENAS}
\author{}
\date{}

\hypersetup{
  pdftitle={Vanguarda 44: Ardenas -- Mini-Livro de Expansão},
  pdfauthor={Felipe Ramires Terrazas},
}

\fancyhead[L]{\small\color{cinzachumbo}Vanguarda 44: Ardenas (Mini-Livro de Expansão)}

\begin{document}

\begin{titlepage}
\centering
\vspace*{3cm}
{\Huge\bfseries\color{vinho} VANGUARDA 44: ARDENAS \par}
\vspace{0.5cm}
{\Large\color{mostarda} Mini-Livro de Expansão \par}
\vspace{2cm}
{\large Cenários Históricos da Batalha das Ardenas \par}
\vfill
\end{titlepage}

\tableofcontents

\bigskip

Dezembro de 1944. A última grande ofensiva alemã no Oeste tentou virar o jogo numa aposta desesperada através da floresta gelada das Ardenas. Este mini-livro traz três cenários ambientados nessa batalha: da Bastogne cercada ao combate lento e sangrento da Floresta de Hürtgen.

\section{Cenários}

\subsection{Bastogne}

\textbf{Contexto Histórico}: Cercada pelo XLVII Panzer Korps alemão (a Divisão Panzer Lehr, a 2ª Divisão Panzer e a 26ª Divisão Volksgrenadier), a 101ª Divisão Aerotransportada defendeu Bastogne por seis dias em pleno inverno, sem suprimentos adequados, até ser aliviada pela 4ª Divisão Blindada do General Patton. Quando intimado a se render em 22 de dezembro, o General Anthony McAuliffe respondeu com uma única palavra que se tornaria célebre: ``Nuts!''

\textbf{Data e Local}: 20 a 26 de dezembro de 1944, Bastogne, Bélgica.

\textbf{Forças}: Alemães (Atacante) vs. Americanos (Defensor, cercado).

\textbf{Terreno}: floresta densa (Terreno Difícil) circundando um vilarejo central (Cobertura Pesada); campo aberto coberto de neve.

\textbf{Implantação}: Americanos ocupam um perímetro ao redor do centro da mesa; Alemães entram de múltiplas bordas simultaneamente, representando o cerco.

\textbf{Qualidade de Tropa}: Americanos Veteranos (mas desgastados); Alemães Regulares.

\textbf{Regra Especial ``Cercados e Sem Suprimentos''}: Americanos só podem usar a ordem Recuperar através do Sargento de cada esquadrão: Tenentes e Capitães não podem emprestar PC extra para Recuperar, refletindo a escassez severa de munição e suprimentos médicos durante o cerco.

\textbf{Objetivo}: Americanos vencem se mantiverem o controle do centro da mesa até o turno 10. Alemães vencem se o tomarem antes disso.

\textbf{Duração}: 10 turnos.

\subsection{Aachen}

\textbf{Contexto Histórico}: Aachen foi a primeira cidade alemã capturada pelos Aliados, num combate urbano brutal travado casa a casa em outubro de 1944. A 1ª Divisão de Infantaria americana enfrentou uma guarnição alemã com ordens de defender a cidade até o último homem, por seu peso simbólico como solo alemão.

\textbf{Data e Local}: outubro de 1944, Aachen, Alemanha.

\textbf{Forças}: Americanos (Atacante) vs. Alemães (Defensor).

\textbf{Terreno}: urbano denso, quase toda a mesa em Cobertura Pesada, com poucas ruas abertas (zonas de matança).

\textbf{Implantação}: Americanos avançam de uma borda; Alemães defendem em profundidade pela cidade.

\textbf{Qualidade de Tropa}: ambos Regulares.

\textbf{Regra Especial ``Defesa até o Último Homem''}: unidades alemãs nunca entram em pânico (Nível 6 de Stress) enquanto houver ao menos 1 outra unidade alemã viva na mesa: continuam lutando, sofrendo as penalidades do Nível 6, até serem a última unidade restante.

\textbf{Objetivo}: Americanos vencem se eliminarem todas as unidades alemãs ou controlarem 75\% da mesa até o turno 10.

\textbf{Duração}: 10 turnos.

\subsection{Floresta de Hürtgen}

\textbf{Contexto Histórico}: Um dos combates mais longos e custosos da campanha americana na Europa foi travado numa floresta densa perto da fronteira alemã, onde a vegetação fechada anulava boa parte da vantagem aérea e blindada americana, forçando um combate de infantaria lento e sangrento entre as árvores, por meses a fio.

\textbf{Data e Local}: setembro a dezembro de 1944, Floresta de Hürtgen, Alemanha.

\textbf{Forças}: Simétrico (Batalha de Encontro), Americanos vs. Alemães, ambos avançando.

\textbf{Terreno}: a mesa inteira é Terreno Difícil (floresta densa), exceto por 1 a 2 pequenas clareiras.

\textbf{Implantação}: bordas opostas da mesa, como uma Batalha de Encontro padrão.

\textbf{Qualidade de Tropa}: ambos Regulares.

\textbf{Regra Especial ``Visibilidade Reduzida''}: o alcance máximo de qualquer arma de tiro direto (Rifle, SMG, LMG) é reduzido pela metade nesta mesa. Morteiros e Snipers não são afetados.

\textbf{Objetivo}: destruir a força inimiga ou controlar o centro da mesa até o turno 8.

\textbf{Duração}: 8 turnos.

\end{document}
